\documentclass[conference]{IEEEtran}
\IEEEoverridecommandlockouts

\usepackage{amsmath}
\usepackage{amssymb}
\usepackage{booktabs}
\usepackage{array}
\usepackage{hyperref}
\usepackage[numbers,sort&compress]{natbib}
\usepackage{enumitem}

\def\BibTeX{{\rm B\kern-.05em{\sc i\kern-.025em b}\kern-.08em
    T\kern-.1667em\lower.7ex\hbox{E}\kern-.125emX}}

\title{Composite Priority Dispatching Rules for the\\
Dynamic Flexible Job-Shop Scheduling Action Space:\\
Selection Rationale and Literature Grounding}
\author{}
\date{\today}

\newcommand{\argmin}{\operatorname*{arg\,min}}
\newcommand{\argmax}{\operatorname*{arg\,max}}

\begin{document}
\maketitle

\begin{abstract}
This note documents the priority dispatching rule (PDR) action space used by the
deep reinforcement learning (DRL) agent in our dynamic flexible job-shop scheduling
(DFJSP) simulator, and the literature basis for each rule. Rather than have the agent
choose directly among all eligible (job, machine) assignments, the agent selects at
each decision point among a small set of \emph{composite} PDRs, each pairing a
job-selection heuristic with a machine-selection (routing) heuristic. We define each
component rule formally, situate it against the classical dispatching-rule literature,
and give the design rationale for the eight-rule set actually used. A short
implementation note at the end records two corrections made to the rule set during a
code audit, for reproducibility. This is an abbreviated working draft, intended as a
basis for the corresponding section of the thesis/paper.
\end{abstract}

\section{Background: Composite Priority Dispatching Rules}
\label{sec:background}

A priority dispatching rule (PDR) is a heuristic scoring function used to rank
candidates at a scheduling decision point, so that the highest (or lowest) scoring
candidate is selected without exhaustive search. PDRs are among the oldest and most
widely deployed scheduling heuristics precisely because they are $O(n)$ per decision
and require no lookahead; \cite{panwalkar1977} catalog and classify over one hundred
such rules for the classical job shop, and \cite{blackstone1982} give a complementary
survey specifically for manufacturing job-shop dispatch. \cite{holthaus1997} and
\cite{rajendran1999} extend the classical PDR literature to the \emph{dynamic} job
shop -- jobs arrive stochastically over time rather than all being known up front,
which is also the setting of our simulator -- and show that combining a job's own
processing-time/work-content information with a measure of downstream congestion
outperforms either signal alone on mean flow time and tardiness.

The flexible job-shop scheduling problem (FJSP) adds a second decision dimension on
top of the classical job shop: each operation may be processed on more than one
eligible machine, so a \emph{routing} (machine-assignment) decision must be made in
addition to the classical \emph{sequencing} (job-priority) decision
\citep{chaudhry2016, xie2019}. Consequently, a PDR for FJSP is naturally expressed as a
pair: a job-selection rule that decides \emph{which} ready job goes next, and a
machine-selection rule that decides \emph{which} eligible machine it goes to. We refer
to such a pair as a \emph{composite} PDR, following the same design pattern used for
the DRL action space in \cite{luo2020}, where a DQN agent selects among six such
composite rules at each decision point rather than acting on raw state directly.

\subsection{Notation}
\label{sec:notation}

Let $\mathcal{J}$ be the set of jobs and $\mathcal{M}$ the set of machines. At
decision time $t$:
\begin{itemize}[leftmargin=*]
  \item For job $j$ and machine $m$, $p_{j,m}$ is the processing time of $j$'s
        current operation on $m$.
  \item $r_j = \sum_{o \geq o_j} \min_{m \in E_o} p_{o,m}$ is job $j$'s estimated
        remaining work: the sum, over $j$'s current and remaining operations, of the
        minimum processing time across each operation's eligible machines.
  \item $a_j$ is job $j$'s arrival time into the system.
  \item $Q_m(t)$ is machine $m$'s queued workload: the sum of processing times of
        all jobs currently queued at $m$ or already committed to it (in transit or
        awaiting pickup), evaluated at $t$.
  \item $U_m(t) = B_m(t) \big/ \big(t - D_m(t)\big)$ is machine $m$'s cumulative
        utilization ratio, where $B_m(t)$ is total busy time accumulated on $m$ up to
        $t$ and $D_m(t)$ is total downtime accumulated on $m$ up to $t$.
\end{itemize}
$Q_m(t)$ and $U_m(t)$ are both congestion signals but operate on different
timescales: $Q_m(t)$ is an instantaneous snapshot of work already committed to $m$,
while $U_m(t)$ is a running average over the whole episode so far. A machine can be
idle right now ($Q_m(t) = 0$) while having run hot all episode ($U_m(t)$ high), or
vice versa.

\section{Job-Selection Rules}
\label{sec:job-rules}

Given the set of jobs $\mathcal{Q}$ ready for a dispatch or routing decision, each
job-selection rule below defines a score to optimize over $\mathcal{Q}$.

\begin{description}[style=nextline, leftmargin=1em]
  \item[SPT (Shortest Processing Time)] $j^\ast = \argmin_{j \in \mathcal{Q}} p_{j,m}$.
    The canonical PDR for minimizing work-in-process and mean flow time
    \citep{panwalkar1977, blackstone1982}; its known weakness is starving long jobs.
  \item[LPT (Longest Processing Time)] $j^\ast = \argmax_{j \in \mathcal{Q}} p_{j,m}$.
    The complementary rule: prioritizing long operations reduces machine-idle
    fragmentation and front-loads throughput risk instead of pushing it to the tail
    \citep{panwalkar1977}.
  \item[SRT (Shortest Remaining Time, i.e., least work remaining)]
    $j^\ast = \argmin_{j \in \mathcal{Q}} r_j$. Prioritizes jobs closest to
    completion, clearing them out of the system quickly; a global (whole-job) analog
    of SPT's local (single-operation) view \citep{holthaus1997}.
  \item[LRT (Longest Remaining Time, i.e., most work remaining)]
    $j^\ast = \argmax_{j \in \mathcal{Q}} r_j$. Front-loads jobs with the most work
    left so they are not squeezed into an increasingly congested tail as the episode
    progresses \citep{holthaus1997, rajendran1999}.
  \item[FIFO (First-Come-First-Served)]
    $j^\ast = \argmax_{j \in \mathcal{Q}} (t - a_j)$, i.e., oldest arrival first. The
    classical fairness/starvation-avoidance rule \citep{panwalkar1977}; unlike the
    other four, its score is based on arrival order rather than work content, making
    it the one job-selection rule in this set that is orthogonal to processing time
    entirely.
\end{description}

\section{Machine-Selection (Routing) Rules}
\label{sec:machine-rules}

Given the set of eligible, available candidate machines $\mathcal{C}$ for a job being
routed, each machine-selection rule below defines a score to optimize over
$\mathcal{C}$.

\begin{description}[style=nextline, leftmargin=1em]
  \item[SMPT (Shortest Machine Processing Time)]
    $m^\ast = \argmin_{m \in \mathcal{C}} p_{j,m}$. The greedy local rule: send the
    job to whichever eligible machine would finish it fastest, ignoring that
    machine's congestion state entirely.
  \item[SRWT (Shortest Remaining Work Time, i.e., minimum queued workload)]
    $m^\ast = \argmin_{m \in \mathcal{C}} Q_m(t)$. A congestion-avoidance rule in the
    work-in-next-queue (WINQ) family: \cite{holthaus1997} and
    \cite{rajendran1999} show that combining a job's own processing time with the
    downstream queue's work content (rather than its head-count) consistently beats
    plain SPT/LPT on mean flow time and tardiness.
  \item[MMUR (Minimum Machine Utilization)]
    $m^\ast = \argmin_{m \in \mathcal{C}} U_m(t)$. A workload-balancing routing rule
    from the same family as ``least-utilized-machine'' assignment rules surveyed for
    FJSP by \cite{xie2019} and \cite{chaudhry2016}, but using the whole-episode
    utilization history rather than the instantaneous queue -- a longer-horizon
    congestion signal, deliberately distinct from SRWT.
\end{description}

\section{The Composite Rule Set}
\label{sec:rule-set}

Rather than take the full cross product of five job-selection and three
machine-selection rules, we use a curated subset of eight composite rules, chosen so
that (a) every job-selection rule and every machine-selection rule above is exercised
at least once, and (b) the pairings are not redundant with one another under the
formal definitions in Sections~\ref{sec:job-rules}--\ref{sec:machine-rules}.
Table~\ref{tab:rules} lists the set as implemented.

\begin{table*}[t]
\centering
\small
\caption{The eight composite PDRs in the action space. A ninth action, \textsc{Random},
uniformly re-samples one of the eight above at every decision; it is included for
baseline/exploration sweeps only and is excluded from the trained agent's action
space (fixed at 8 discrete actions).}
\label{tab:rules}
\begin{tabular}{@{}lllp{11cm}@{}}
\toprule
\textbf{Rule} & \textbf{Job selection} & \textbf{Machine selection} & \textbf{Rationale} \\
\midrule
SPT\_SMPT & SPT & SMPT & Pure local-greedy throughput baseline. \\
SPT\_SRWT & SPT & SRWT & Throughput-oriented job choice with congestion-aware routing. \\
LPT\_MMUR & LPT & MMUR & Long-operation priority balanced against long-horizon machine load. \\
LPT\_SMPT & LPT & SMPT & Long-operation priority with purely local machine choice. \\
SRT\_SRWT & SRT & SRWT & Clear near-complete jobs quickly, avoid congested machines. \\
SRT\_SMPT & SRT & SMPT & Clear near-complete jobs quickly, purely local machine choice. \\
LRT\_MMUR & LRT & MMUR & Front-load heavy jobs, balanced against long-horizon machine load. \\
FIFO\_SRWT & FIFO & SRWT & Fairness/arrival-order priority with congestion-aware routing. \\
\bottomrule
\end{tabular}
\end{table*}

By construction, every rule in Table~\ref{tab:rules} has a distinct
(job-selection, machine-selection) pair, so no two rules make identical decisions
given identical inputs -- a property that did not hold for an earlier revision of
this rule set (see Section~\ref{sec:implementation}).

\section{Implementation Note (for reproducibility)}
\label{sec:implementation}

Two corrections were made to the rule set described above during a code audit, after
an earlier revision's results showed the MMUR-paired rules behaving in a way
inconsistent with their intended definition:

\begin{enumerate}[leftmargin=*]
  \item \textbf{MMUR routing direction.} The machine-selection half of MMUR was
        originally implemented as $\argmax_{m \in \mathcal{C}} Q_m(t)$ -- routing
        each job to the \emph{most} congested candidate machine -- rather than the
        $\argmin$ given in Section~\ref{sec:machine-rules}. This produced a
        positive-feedback overload effect (busy machines were preferentially made
        busier), consistent with the near-total non-completion of
        LPT\_MMUR/LRT\_MMUR jobs observed at higher load in earlier sweeps. The rule
        now implements the $\argmin$ form above. Additionally, MMUR's
        machine-selection signal was changed from an instantaneous queue-length
        reading to the cumulative utilization ratio $U_m(t)$ defined in
        Section~\ref{sec:notation}, to give it a genuinely distinct basis from SRWT
        (which already used the instantaneous queued-workload signal).
  \item \textbf{Removal of a redundant rule.} An intermediate revision also included
        a ninth rule, \textsc{MRT\_SRWT} (job selection: LRT's $\argmax_j r_j$;
        machine selection: SRWT), added specifically as an unconfounded baseline
        while MMUR's machine-selection logic was still defective. Once MMUR was
        corrected per item 1, \textsc{MRT\_SRWT} and \textsc{LRT\_MMUR} shared an
        identical job-selection function and differed only in machine-selection
        signal (instantaneous queued workload vs.\ cumulative utilization) -- too
        narrow a distinction to justify two dedicated action-space slots. It was
        removed; \textsc{LRT\_MMUR} is retained as the intended pairing.
\end{enumerate}

All results generated prior to these corrections used the pre-correction (reversed)
MMUR definition and/or included \textsc{MRT\_SRWT}, and should not be interpreted
against the rule definitions given in Sections~\ref{sec:job-rules}--\ref{sec:machine-rules}.

\bibliographystyle{unsrtnat}
\bibliography{references}

\end{document}
